\section{The polar source column has an exact rank-two quotient defect}
\label{sec:polar-quotient-rank}

Under the half-plane hypothesis of MA-MO-018, all nontrivial
zeros lie strictly inside the completed strip. This lets the
finite outgoing condition of MA-MO-016 be decided exactly.
The two polar columns remain independent after passage to the
actual zero quotient, for every nonzero compact test. The
argument retains the full multiplicity of every zero.

Keep \(0\le a_*<a<1/2\), \(s\ge0\), and the explicit
zero-free half-plane (MA-MO.320). Write
\[
 \begin{gathered}
 H=H_{\Sigma,a,s},\qquad J=\overline{\Theta\mathscr U_0}^{\,H},
 \qquad\mathcal W=H\oplus\mathbb C^2,
 \qquad\mathcal W_+=J\oplus\mathbb C^2,\\
 \mathfrak q:\mathcal W\longrightarrow H/J,
                    \qquad\mathfrak q(y,c)=y+J.
 \end{gathered}
 \tag{MA-MO.400}
\]
The source sections are the actual Gaussian vectors
\(s_j=i_+f_j\) of (MA-MO.251), and \(r_j\) are the
actual functions in (MA-MO.255).

\paragraph{MA-MO-020 (exact outgoing rank and noninvariance).}
For every nonzero \(g\in C_c^\infty(\mathbb R)\), the
map
\[
 \begin{aligned}
 \mathbb C^2&\longrightarrow H/J,\\
 (c_0,c_1)&\longmapsto[C_g(c_0r_0+c_1r_1)]
 \end{aligned}
 \tag{MA-MO.401}
\]
is injective. Consequently
\[
        \operatorname{rank}
         (\mathfrak q\mathscr D_g|_{\mathcal W_+})=2,
                \qquad
        \mathscr D_g(\mathcal W_+)\not\subseteq\mathcal W_+.
 \tag{MA-MO.402}
\]
Here \(\mathscr D_g\) is the actual full trace-class matrix
of MA-MO-018, and its positive source restriction is the map
of MA-MO-016. A correction making that source restriction
take values in \(\mathcal W_+\) must have rank at least two.
The orthogonal correction
\[
       (1-P_{\mathcal W_+})
                        \mathscr D_g|_{\mathcal W_+}
 \tag{MA-MO.403}
\]
has rank exactly two and attains that minimum. It changes
the original source action.

The conclusions still hold when \(g(0)=0\), and even when
\(\mathcal Mg(1/2)=\mathcal Mg(-1/2)=0\), provided
\(g\ne0\). They concern this particular completed source
image and the original operator. They are not an obstruction
to every possible return map or relative-trace construction.

\begin{proof}
\emph{The exact two source Mellin residues.}
Poisson summation and the Gaussian Fourier interchange give
\[
       r_0=\Theta f_0+\phi e^{x/2},\qquad
       r_1=\Theta f_1-\phi e^{-x/2}.
 \tag{MA-MO.404}
\]
For example \(R\Theta f_1=\Theta f_0+e^{x/2}\),
whereas \(R\Theta f_0=\Theta f_1-e^{-x/2}\).
Substitution in (MA-MO.255) proves both signs in
(MA-MO.404). Each \(r_j\) lies in \(\mathcal E\):
the displayed term cancels its left Poisson tail, and the
cutoff vanishes at the right.

Let
\[
 \begin{aligned}
 F_j(x)&=2e^{x/2}f_j(e^x),&Z(z)&=\zeta(1/2+z),\\
 P(z)&=\int e^{zx}\phi'(x)\,dx,&P(0)&=-1.
 \end{aligned}
 \tag{MA-MO.405}
\]
The Gaussian integral and gamma recurrence give, first for
\(\Re z>-1/2\),
\[
 \mathcal MF_0(z)=(1/2-z)\Gamma_{\mathbb R}(1/2+z),
 \qquad
 \mathcal MF_1(z)=(1/2+z)\Gamma_{\mathbb R}(1/2+z).
 \tag{MA-MO.406}
\]
In particular both are holomorphic in the open strip
\(|\Re z|<1/2\). The second has value one at \(z=1/2\),
and the first has value zero there, in agreement with the
additive moments of the sections.

On \(\Re z>1/2\), absolute theta summation gives
\(\mathcal M\Theta f_j(z)=Z(z)\mathcal MF_j(z)\).
Integration by parts gives
\(\int e^{\xi x}\phi(x)\,dx=-P(\xi)/\xi\)
when \(\Re\xi>0\). Thus (MA-MO.404) gives the exact
continued identities
\[
 \begin{aligned}
 \mathcal Mr_0(z)
     &=Z(z)\mathcal MF_0(z)-\frac{P(z+1/2)}{z+1/2},\\
 \mathcal Mr_1(z)
     &=Z(z)\mathcal MF_1(z)+\frac{P(z-1/2)}{z-1/2}.
 \end{aligned}
 \tag{MA-MO.407}
\]
They hold throughout the critical open strip by meromorphic
continuation. The left sides are entire. At \(z=1/2\)
the pole in the second formula cancels because its zeta
residue is one and \(P(0)=-1\). At \(z=-1/2\) the
apparent pole in the first formula cancels its Gaussian
term: \(\Gamma_{\mathbb R}(w)=2/w+O(1)\) and
\(\zeta(0)=-1/2\) give residue \(-1\) in its first
term and residue \(+1\) in its second. The value
\(\zeta(0)=-1/2\) also follows directly from the reflection
formula and the residue one at \(s=1\). Neither endpoint
is a nontrivial zero in the argument
below, so no endpoint division is needed.

If \(\lambda=\rho-1/2\) is a nontrivial zero of
multiplicity \(m_\rho\), the holomorphic first terms in
(MA-MO.407) vanish to order at least \(m_\rho\) there.
For \(c=(c_0,c_1)\), put
\[
 \begin{aligned}
 E_c(z)&=-c_0(z-1/2)P(z+1/2)
                       +c_1(z+1/2)P(z-1/2),\\
 E_c(1/2)&=-c_1,&E_c(-1/2)&=-c_0.
 \end{aligned}
 \tag{MA-MO.408}
\]
Thus \(E_c\) is a nonzero entire function whenever
\(c\ne0\). These two evaluations are legitimate values
of the entire numerator; they are not evaluations of the
singular fractions in (MA-MO.407).

\emph{Full retained jets, not only point evaluations.}
The functional equation reflects every nontrivial zero
\(\rho\) to \(1-\overline\rho\). Hypothesis (MA-MO.320)
therefore implies
\[
          |\Re(\rho-1/2)|\le a_*<a.
 \tag{MA-MO.409}
\]
By the exact sum-weight closure theorem MA-WP-001, every
zero is strictly interior and all of its
\(m_\rho\) jets are retained, for every \(s\ge0\).
Suppose \(C_g(c_0r_0+c_1r_1)\in J\), and write
\(G(z)=\mathcal Mg(z)\). Its Mellin transform is
\(G(z)\mathcal M(c_0r_0+c_1r_1)(z)\).
It must vanish to order at least \(m_\rho\) at every
\(\lambda=\rho-1/2\). Equations (MA-MO.407)--(MA-MO.408)
then show that
\[
                G(z)E_c(z)
   \text{ vanishes to order at least }m_\rho
                     \text{ at every }z=\rho-1/2.
 \tag{MA-MO.410}
\]
Indeed multiplication by \(z^2-1/4\) is invertible in
each of these local jet rings. No simplicity assumption
has entered this implication.

\emph{The exponential-type zero count.}
If \(c\ne0\) and \(g\ne0\), the product \(GE_c\)
is not identically zero, by Fourier uniqueness for \(g\)
and (MA-MO.408). Compactness of the supports of \(g\)
and \(\phi'\) gives constants \(C,R\) such that
\[
             |G(z)E_c(z)|\le C(1+|z|)e^{R|z|}
                                     \quad(z\in\mathbb C).
 \tag{MA-MO.411}
\]
Here the factors \(e^{\pm x/2}\) on the fixed support
of \(\phi'\) only change \(C\).

To spell out the zero-count consequence, choose a point
\(z_0\) where \(G(z_0)E_c(z_0)\ne0\). Jensen's formula
on the disk of radius \(2r\) about \(z_0\) gives
\[
 n(r)\log2
 \le\frac1{2\pi}\int_0^{2\pi}
       \log|G(z_0+2re^{it})E_c(z_0+2re^{it})|\,dt
           -\log|G(z_0)E_c(z_0)|=O(r).
 \tag{MA-MO.412}
\]
Here \(n(r)\) counts zeros in the smaller disk, with
multiplicity. Radii meeting a zero are handled by a limit
or a slightly larger radius. Thus (MA-MO.410) would imply
that the full nontrivial zeta divisor has only \(O(T)\)
zeros of height at most \(T\), counted with multiplicity.

\emph{Why the full zeta divisor cannot have that count.}
We prove the needed incompatibility, so a distinct-zero
count or an unexamined asymptotic is not being substituted.
Let
\[
       \xi(s)=\tfrac12s(s-1)\pi^{-s/2}\Gamma(s/2)\zeta(s).
 \tag{MA-MO.413}
\]
It is entire of order at most one, with nonzero values
at zero and one and with precisely the nontrivial zeta
divisor. The order bound follows from the Euler bound on
\(\Re s\ge2\), the functional equation on
\(\Re s\le-1\), and the fixed-strip polynomial zeta
estimate and gamma growth between them, exactly as in the
Hadamard derivation (MA-RB.8)--(MA-RB.10).

Suppose its zero count, including multiplicities for
\(|\Im\rho|\le T\), were \(O(T)\). Hadamard
factorization and logarithmic differentiation give
\(\xi'/\xi(s)=B+\sum_\rho((s-\rho)^{-1}+\rho^{-1})\),
where \(B=\xi'(0)/\xi(0)\) is real. All sums of real
parts below converge absolutely: \(0<\Re\rho<1\)
and the assumed count imply convergence of
\(\sum_{|\rho|>1}|\rho|^{-2}\).
Put \(S=\sum_\rho\Re(1/\rho)\). The functional equation
at zero and one, and symmetry under \(\rho\mapsto1-\rho\),
give \(-B=B+2S\). Therefore for real \(\sigma>1\),
\[
       \frac{\xi'}{\xi}(\sigma)
          =\sum_\rho\frac{\sigma-\Re\rho}
                    {(\sigma-\Re\rho)^2+(\Im\rho)^2}.
 \tag{MA-MO.414}
\]
For \(\sigma\ge2\), each summand is at most
\(\sigma/(\sigma^2/4+(\Im\rho)^2)\). The zeros with
\(|\Im\rho|\le\sigma\) contribute \(O(1)\), using
their assumed \(O(\sigma)\) count. The dyadic shell
\(2^k\sigma<|\Im\rho|\le2^{k+1}\sigma\) contributes
\(O(2^{-k})\). Consequently the whole right side of
(MA-MO.414) is \(O(1)\), uniformly as \(\sigma\to\infty\).

On the other hand, direct differentiation of (MA-MO.413)
gives
\[
 \frac{\xi'}{\xi}(\sigma)
   =\frac1\sigma+\frac1{\sigma-1}-\frac12\log\pi
       +\frac12\psi(\sigma/2)
       +\frac{\zeta'}{\zeta}(\sigma).
 \tag{MA-MO.415}
\]
The final term tends to zero by its absolutely convergent
Euler series, dominated for \(\sigma\ge2\). For \(x>1\),
let \(U\) have density \(u^{x-1}e^{-u}/\Gamma(x)\).
Then \(\psi(x)=\mathbb E\log U\) and
\(\mathbb E(U^{-1})=1/(x-1)\). Jensen applied to
\(\log(U^{-1})\) proves
\(\psi(x)\ge\log(x-1)\). Thus (MA-MO.415) tends to
\(+\infty\), contradicting the bounded value of
(MA-MO.414). This proves the required failure of an
\(O(T)\) full zero count, including every multiplicity.
The contradiction with (MA-MO.412) proves that membership
in \(J\) was impossible for \(c\ne0\).

\emph{The exact outgoing rank and correction.}
We have proved injectivity of (MA-MO.401). Formula
(MA-MO.256) says that the outgoing map of the positive
source column is its negative, applied to the two polar
coordinates, and vanishes on its \(J\) summand. Its
rank is consequently exactly two. If a linear correction
\(C:\mathcal W_+\to\mathcal W\) makes
\(\mathscr D_g-C\) take values in \(\mathcal W_+\),
then \(\mathfrak qC=\mathfrak q\mathscr D_g\) has rank
two, so \(\operatorname{rank}C\ge2\). Orthogonal
projection onto the complement of \(\mathcal W_+\)
identifies that complement with \(H/J\); hence (MA-MO.403)
has rank exactly two. This proves all the assertions.
\end{proof}

\paragraph{Counting input and framework scope.}
The only zero-location use is (MA-MO.409), which makes the
entire nontrivial divisor retained in the actual Hilbert
quotient. The counting contradiction itself is unconditional.
For comparison, the stronger classical count is explicitly
quantified by Hasanalizade, Shen, and Wong (2022),
Corollary~1.2, equation~(1.5), p.~221; its completed
normalization and argument-principle setup are in \S2,
equations~(2.1)--(2.3), pp.~222--223. The argument above
does not require their numerical estimate.

The primary reference is: Hasanalizade, E., Shen, Q., \&
Wong, P.-J. (2022). Counting zeros of the Riemann zeta
function. \emph{Journal of Number Theory, 235}, 219--241.
\url{https://doi.org/10.1016/j.jnt.2021.06.032}.
The completed normalization and reflection identity also
agree with NIST DLMF, equations~25.4.3--25.4.4.
Only the stated source locators and counting setup were
checked here; the numerical estimate is corroborative,
not a new unproved premise of MA-MO-020.

At each retained zero the argument uses the full local
finite-length factor \(\mathbb C[z]/(z^{m_\rho})\),
with all its nilpotent jets. The inherited cote~1 bookkeeping
therefore retains precisely the multiplicities needed in
Jensen and Hadamard. The rank-two orthogonal correction
is a concrete changed analytic map. Neither a finite-length
\(K_0\) identity nor a weight assertion makes the original
source image invariant or assigns a quotient trace to the
uncorrected matrix.
